% ============================================================
%  Artificial Intelligence: A Textbook
%  Chapter 3 — Multiagent Search
%  Offline Beamer slides (uses shared theme.tex + branding.tex)
%  Prince Sattam Bin Abdulaziz University — CCES
% ============================================================
\documentclass[aspectratio=169,11pt]{beamer}

\input{theme.tex}

% ---------- COURSE RE-BRANDING (AI instead of CS3401) --------
\graphicspath{{../chapters/images/ch3/}{../chapters/images/ch4/}{../chapters/images/}}
% ============================================================
%  Artificial Intelligence: A Textbook
%  Shared re-branding for the AI course (input AFTER theme.tex)
%  Overrides the CS3401 footer/title page from theme.tex.
% ============================================================

% Footer rebranded for the AI course
\setbeamertemplate{footline}{%
  \begin{beamercolorbox}[wd=\paperwidth,ht=0.4pt,dp=0pt]{footer rule}{}
  \end{beamercolorbox}
  \leavevmode
  \hbox{%
    \begin{beamercolorbox}[wd=0.60\paperwidth,ht=1.8em,dp=0.6em,leftskip=1em]{normal text}%
      \textcolor{PrimaryDark}{\tiny\textbf{Dr. Khalil Chebil} %
        \quad\textbar\quad CCES, PSAU}%
    \end{beamercolorbox}%
    \begin{beamercolorbox}[wd=0.40\paperwidth,ht=1.8em,dp=0.6em,rightskip=1em,right]{normal text}%
      \textcolor{PrimaryDark}{\tiny\insertshorttitle%
        \quad\textbar\quad\insertframenumber\,/\,\inserttotalframenumber}%
    \end{beamercolorbox}%
  }%
}

% Title page rebranded for the AI course
\setbeamertemplate{title page}{%
  \begin{tikzpicture}[remember picture,overlay]
    \fill[PrimaryDark](current page.north west) rectangle (current page.south east);
    \fill[PrimaryMid,opacity=0.45]
      ([shift={(2.5cm,-2.5cm)}]current page.north east) circle(5.5cm);
    \fill[AccentOrange,opacity=0.10]
      ([shift={(-1.5cm,2.0cm)}]current page.south west) circle(4.2cm);
    \draw[white,opacity=0.08,line width=40pt]
      ([shift={(0cm,-0.5cm)}]current page.north west)
      to[out=-30,in=210]
      ([shift={(0cm,0.5cm)}]current page.south east);
  \end{tikzpicture}%
  \vspace{0.6cm}
  \begin{center}
    {\tiny\textcolor{PrimaryLight}{\textbf{CCES} \textendash{}
      College of Computer Engineering and Sciences}}\\[0.2em]
    {\tiny\textcolor{PrimaryLight}{Prince Sattam Bin Abdulaziz University}}\\[1.2em]
    {\footnotesize\textcolor{AccentYellow}{\textbf{CS600: Advanced AI }}}\\[1.5em]
    \textcolor{PrimaryMid}{\rule{0.55\textwidth}{0.6pt}}\\[0.9em]
    {\Large\usebeamerfont{title}\textcolor{white}{\inserttitle}}\\[0.4em]
    {\small\usebeamerfont{subtitle}\textcolor{PrimaryLight}{\insertsubtitle}}\\[0.7em]
    \textcolor{PrimaryMid}{\rule{0.55\textwidth}{0.6pt}}\\[1.0em]
    {\small\textcolor{white}{\insertauthor}}\\[0.3em]
    {\footnotesize\textcolor{PrimaryLight}{\insertdate}}
  \end{center}%
}


% ---------- SLIDE METADATA ----------------------------------
\title[Ch.\,3 — Multiagent Search]{Multiagent Search}
\subtitle{Chapter 3 — Adversarial Games: Minimax, $\alpha$–$\beta$, and MCTS}
\author{Dr.\ Khalil Chebil}
\date{Artificial Intelligence Course}

% Inline check / cross marks
\newcommand{\cmark}{\textcolor{SuccessGreen}{\ensuremath{\checkmark}}}
\newcommand{\xmark}{\textcolor{AlertRed}{\ensuremath{\times}}}

\begin{document}

% ============================================================
%  TITLE
% ============================================================
{\setbeamertemplate{footline}{}
\begin{frame}[plain,noframenumbering]
  \titlepage
\end{frame}}

% ============================================================
%  OUTLINE
% ============================================================
\begin{frame}{Chapter Outline}{What we will cover}
  \begin{columns}[T]
    \begin{column}{0.5\textwidth}
      \tableofcontents[sections={1-3}]
    \end{column}
    \begin{column}{0.5\textwidth}
      \tableofcontents[sections={4-6}]
    \end{column}
  \end{columns}
  \vfill
  \begin{block}{Learning Objectives}
    \begin{itemize}
      \item Understand \keyword{multiagent} environments and strategy selection
      \item Model contingencies with \keyword{AND-OR} search trees
      \item Apply \keyword{minimax} and \keyword{alpha--beta} pruning to adversarial games
      \item Use \keyword{evaluation functions} and \keyword{Monte Carlo Tree Search}
    \end{itemize}
  \end{block}
\end{frame}

% ============================================================
\section{Multiagent Environments}
% ============================================================

\begin{frame}{From Plan Finding to Strategy Selection}{Others react to your moves}
  \begin{columns}[T]
    \begin{column}{0.55\textwidth}
      In \keyword{single-agent} search, an action determines the next state.
      In \keyword{multiagent} search, other agents cooperate, compete, or act
      independently.
      \vspace{0.3em}
      \begin{alertblock}{The central shift}
        \small A move does not define \emph{one} future — it defines a
        \keyword{set of possible futures}. We seek a \keyword{strategy}, not a
        fixed plan.
      \end{alertblock}
    \end{column}
    \begin{column}{0.45\textwidth}
      \begin{block}{Key challenges}
        \small
        \begin{itemize}
          \item Others have different goals \& information
          \item Effects depend on others' responses
          \item Needs contingency reasoning
          \item Search grows fast with each reply
        \end{itemize}
      \end{block}
    \end{column}
  \end{columns}
\end{frame}

\begin{frame}{Types of Multiagent Environments}{Method depends on the relationship}
  \begin{columns}[T]
    \begin{column}{0.34\textwidth}
      \begin{block}{Competitive}
        \small Conflicting objectives; one's gain is another's loss.
        \\[0.3em] \scriptsize \textit{Chess, Go, Tic-Tac-Toe} \\ $\Rightarrow$ adversarial search
      \end{block}
    \end{column}
    \begin{column}{0.34\textwidth}
      \begin{exampleblock}{Independent}
        \small Act without explicit coordination, yet still affect one another.
        \\[0.3em] \scriptsize \textit{Self-driving cars in traffic}
      \end{exampleblock}
    \end{column}
    \begin{column}{0.34\textwidth}
      \begin{block}{Cooperative}
        \small Share information or align objectives for a joint outcome.
        \\[0.3em] \scriptsize \textit{Warehouse robots, rescue teams}
      \end{block}
    \end{column}
  \end{columns}
  \vspace{0.5em}
  \begin{center}\small
    \important{Tic-Tac-Toe:} after you place \texttt{X}, the opponent chooses
    \texttt{O} — your move opens many futures, not one.
  \end{center}
\end{frame}

% ============================================================
\section{AND-OR Search Trees}
% ============================================================

\begin{frame}{AND-OR Search Trees}{Our choices vs.\ outcomes we must handle}
  \begin{columns}[c]
    \begin{column}{0.5\textwidth}
      \begin{itemize}
        \item \keyword{OR node} — \emph{our} turn: pick \keyword{one} good action
        \item \keyword{AND node} — outside our control: the plan must work for
          \keyword{all} relevant continuations
      \end{itemize}
      \vspace{0.3em}
      \begin{block}{In adversarial games}
        \small OR $=$ our move \quad AND $=$ opponent's possible replies.
        \\[0.2em]\scriptsize (AND nodes may also model uncertain environment outcomes.)
      \end{block}
    \end{column}
    \begin{column}{0.5\textwidth}
      \begin{center}
      \begin{tikzpicture}[level distance=1.05cm, node distance=0.5cm,
          orn/.style={circle, draw=PrimaryMid, fill=PrimaryLight!50, minimum size=0.7cm, font=\scriptsize\bfseries},
          andn/.style={rectangle, rounded corners=1pt, draw=AccentOrange, fill=AccentYellow!35, minimum size=0.6cm, font=\scriptsize\bfseries},
          level 1/.style={sibling distance=2.6cm},
          level 2/.style={sibling distance=1.3cm}]
        \node[orn] {OR}
          child {node[andn] {AND}
            child {node[orn] {OR}}
            child {node[orn] {OR}}}
          child {node[andn] {AND}
            child {node[orn] {OR}}};
      \end{tikzpicture}\\[0.3em]
      \scriptsize One good choice at OR; ready for every branch at AND.
      \end{center}
    \end{column}
  \end{columns}
\end{frame}

\begin{frame}{AO-Search \& Many Agents}{Guaranteed paths in an AND-OR tree}
  \begin{columns}[T]
    \begin{column}{0.5\textwidth}
      \begin{block}{AO-search logic}
        \small
        \begin{itemize}
          \item At an \keyword{OR} node: succeed if \keyword{some} action succeeds
          \item At an \keyword{AND} node: succeed only if \keyword{every} outcome succeeds
        \end{itemize}
      \end{block}
      \centering\scriptsize A plan is a \emph{sub-tree} guaranteeing the goal.
    \end{column}
    \begin{column}{0.5\textwidth}
      \begin{exampleblock}{More than two agents}
        \small In a 3-agent game (A, B, C):
        \begin{itemize}
          \item OR nodes $=$ agent A's choices
          \item AND nodes $=$ combined choices of B and C
        \end{itemize}
      \end{exampleblock}
      \centering\scriptsize \important{Combining actions blows up the branching factor.}
    \end{column}
  \end{columns}
\end{frame}

\begin{frame}{The AND-OR Tree, Illustrated}{}
  \begin{center}
    \includegraphics[width=0.6\textwidth]{1.png}
  \end{center}
  \begin{center}\scriptsize
    Our agent chooses at OR nodes; the plan must stay valid across the outcomes at AND nodes.
  \end{center}
\end{frame}

% ============================================================
\section{Evaluation \& Depth-Limited Search}
% ============================================================

\begin{frame}{Evaluation Functions}{``How promising is this state right now?''}
  \begin{columns}[T]
    \begin{column}{0.55\textwidth}
      Full trees are infeasible in chess or Go. An \keyword{evaluation function}
      estimates the quality of a \keyword{non-terminal} state.
      \vspace{0.3em}
      \begin{block}{Chess features}
        \small material advantage \textbullet\ king safety \textbullet\ center
        control \textbullet\ piece activity
      \end{block}
    \end{column}
    \begin{column}{0.45\textwidth}
      \begin{alertblock}{Powerful but imperfect}
        \small
        \begin{itemize}
          \item Summarize rich positions with few features
          \item May miss deep tactical ideas
          \item Their quality drives decision quality
        \end{itemize}
      \end{alertblock}
    \end{column}
  \end{columns}
\end{frame}

\begin{frame}{Terminology \& Depth-Limited Search}{Fix the vocabulary}
  \begin{columns}[T]
    \begin{column}{0.5\textwidth}
      \begin{block}{Core terms}
        \small
        \begin{itemize}
          \item \keyword{State / Action / Terminal state}
          \item \keyword{Utility}: larger is better for MAX
          \item \keyword{Loss / cost}: smaller is better
          \item \keyword{MAX / MIN}: the two adversarial roles
        \end{itemize}
      \end{block}
    \end{column}
    \begin{column}{0.5\textwidth}
      \begin{exampleblock}{Depth-limited search}
        \small Cut off expansion at depth $d$, then apply the evaluation function
        at the frontier.
        \begin{itemize}
          \item Deeper search $\Rightarrow$ better decisions
          \item Too deep $\Rightarrow$ infeasible
          \item Good evaluation offsets shallow search
        \end{itemize}
      \end{exampleblock}
    \end{column}
  \end{columns}
  \vspace{0.2em}
  \begin{center}\scriptsize
    Strong engines combine \keyword{lookahead} with a good \keyword{evaluation function}.
  \end{center}
\end{frame}

% ============================================================
\section{Minimax}
% ============================================================

\begin{frame}{Minimax for Adversarial Search}{Optimal play in zero-sum games}
  \begin{columns}[T]
    \begin{column}{0.5\textwidth}
      For two-player, \keyword{zero-sum}, deterministic, perfect-information games.
      \begin{itemize}\small
        \item \keyword{MAX} maximizes the utility value
        \item \keyword{MIN} minimizes it
        \item Terminal states give the true payoff
        \item Others inherit the value of the best move for the player to act
      \end{itemize}
    \end{column}
    \begin{column}{0.5\textwidth}
      \begin{block}{Mutual recursion}
        \scriptsize
        \texttt{max\_player(s)} $=\displaystyle\max_{a}$ \texttt{min\_player(s')}\\[0.2em]
        \texttt{min\_player(s)} $=\displaystyle\min_{a}$ \texttt{max\_player(s')}\\[0.2em]
        base case: depth $=0$ or terminal $\Rightarrow$ \texttt{utility(s)}
      \end{block}
      \centering\scriptsize A model of \important{perfect play} — the best
      \emph{guaranteed} outcome.
    \end{column}
  \end{columns}
\end{frame}

\begin{frame}{Minimax Example}{Read the tree from the leaves upward}
  \begin{center}
  \begin{tikzpicture}[level distance=1.25cm,
      maxn/.style={rectangle, rounded corners=2pt, draw=PrimaryDark, fill=PrimaryLight!55, font=\small\bfseries, minimum size=0.7cm},
      minn/.style={diamond, draw=AlertRed, fill=AlertRed!15, font=\small\bfseries, inner sep=1pt, minimum size=0.7cm},
      leaf/.style={circle, draw=PrimaryMid, fill=white, font=\scriptsize, minimum size=0.55cm},
      level 1/.style={sibling distance=4.2cm},
      level 2/.style={sibling distance=1.35cm},
      edge from parent/.style={draw=PrimaryMid}]
    \node[maxn] (root) {MAX $=3$}
      child {node[minn] {$3$}
        child {node[leaf] {3}}
        child {node[leaf] {12}}
        child {node[leaf] {8}}}
      child {node[minn] {$2$}
        child {node[leaf] {2}}
        child {node[leaf] {4}}
        child {node[leaf] {6}}}
      child {node[minn] {$2$}
        child {node[leaf] {14}}
        child {node[leaf] {5}}
        child {node[leaf] {2}}};
  \end{tikzpicture}
  \end{center}
  \begin{center}\small
    MIN nodes: $\min(3,12,8){=}3$, $\min(2,4,6){=}2$, $\min(14,5,2){=}2$ \quad
    $\Rightarrow$ \quad MAX: \textcolor{AlertRed}{$\max(3,2,2)=\mathbf{3}$} (left branch).
  \end{center}
\end{frame}

% ============================================================
\section{Alpha--Beta \& MCTS}
% ============================================================

\begin{frame}{Alpha--Beta Pruning}{Same answer as minimax, fewer nodes}
  \begin{columns}[T]
    \begin{column}{0.5\textwidth}
      Skip branches that cannot change the final choice.
      \begin{itemize}\small
        \item $\alpha$ — best value \keyword{MAX} can guarantee so far
        \item $\beta$ — best value \keyword{MIN} can guarantee so far
        \item Prune the branch when $\alpha \ge \beta$
      \end{itemize}
      \begin{alertblock}{Move ordering matters}
        \small Examining strong moves early creates tighter bounds and prunes more.
      \end{alertblock}
    \end{column}
    \begin{column}{0.5\textwidth}
      \begin{center}
      \begin{tikzpicture}[level distance=1.15cm,
          maxn/.style={rectangle, rounded corners=2pt, draw=PrimaryDark, fill=PrimaryLight!55, font=\scriptsize\bfseries, minimum size=0.6cm},
          minn/.style={diamond, draw=AlertRed, fill=AlertRed!15, font=\scriptsize\bfseries, inner sep=1pt, minimum size=0.6cm},
          leaf/.style={circle, draw=PrimaryMid, fill=white, font=\scriptsize, minimum size=0.5cm},
          pruned/.style={circle, draw=gray, fill=gray!20, font=\scriptsize, text=gray, minimum size=0.5cm},
          level 1/.style={sibling distance=2.9cm},
          level 2/.style={sibling distance=0.85cm},
          edge from parent/.style={draw=PrimaryMid}]
        \node[maxn] {MAX}
          child {node[minn] {$3$}
            child {node[leaf] {3}}
            child {node[leaf] {12}}
            child {node[leaf] {8}}}
          child {node[minn] {$\le 2$}
            child {node[leaf] {2}}
            child {node[pruned] {4}}
            child {node[pruned] {6}}};
      \end{tikzpicture}\\[0.2em]
      \scriptsize After leaf $2$, the right MIN $\le 2 < 3=\alpha$ — prune the rest.
      \end{center}
    \end{column}
  \end{columns}
\end{frame}

\begin{frame}{Alpha--Beta: Worked Example}{Track $\alpha$ and $\beta$ after each child}
  \begin{columns}[c]
    \begin{column}{0.5\textwidth}
      \centering
      \includegraphics[width=\linewidth]{3.png}\\[0.2em]
      \scriptsize Where does pruning occur (left-to-right)?
    \end{column}
    \begin{column}{0.5\textwidth}
      \centering
      \includegraphics[width=\linewidth]{4.png}\\[0.2em]
      \scriptsize Solution: a node already worse than an earlier alternative
      cannot improve the parent's decision.
    \end{column}
  \end{columns}
\end{frame}

\begin{frame}{Monte Carlo Tree Search (MCTS)}{Estimate by simulation, not exhaustion}
  \begin{columns}[T]
    \begin{column}{0.52\textwidth}
      Ideal for \keyword{huge branching factors} or when a good evaluation
      function is hard to design.
      \vspace{0.2em}
      \begin{block}{Four phases}
        \small
        \begin{enumerate}
          \item \keyword{Selection} — descend via UCT
          \item \keyword{Expansion} — add a child at the frontier
          \item \keyword{Simulation} — roll out to estimate quality
          \item \keyword{Backpropagation} — update visits \& rewards
        \end{enumerate}
      \end{block}
      \centering\scriptsize \emph{Anytime}: more simulations $\Rightarrow$ better estimates.
    \end{column}
    \begin{column}{0.48\textwidth}
      \begin{exampleblock}{UCT: exploit vs.\ explore}
        \[ u_i = \frac{w_i}{n_i} + c\,\sqrt{\frac{\ln N_i}{n_i}} \]
        \scriptsize
        $w_i$ reward, $n_i$ visits to $i$, $N_i$ parent visits, $c$ exploration weight.
      \end{exampleblock}
      \begin{center}
      \begin{tikzpicture}[node distance=0.3cm,
          p/.style={processstep, minimum width=3.2cm, minimum height=0.42cm, font=\scriptsize}]
        \node[p, fill=PrimaryLight!45] (s) {Select};
        \node[p, right=0.3cm of s, fill=SuccessGreen!20] (e) {Expand};
        \node[p, below=0.32cm of s, fill=AccentYellow!35] (sim) {Simulate};
        \node[p, right=0.3cm of sim, fill=AlertRed!15] (b) {Backprop.};
        \draw[-{Latex[length=1.4mm]},PrimaryMid] (s)--(e);
        \draw[-{Latex[length=1.4mm]},PrimaryMid] (e.south) to[out=-90,in=90] (b.north);
        \draw[-{Latex[length=1.4mm]},PrimaryMid] (b)--(sim);
        \draw[-{Latex[length=1.4mm]},PrimaryMid] (sim.north) to[out=90,in=-90] (s.south);
      \end{tikzpicture}
      \end{center}
    \end{column}
  \end{columns}
\end{frame}

% ============================================================
\section{Choosing a Method \& Summary}
% ============================================================

\begin{frame}{Choosing the Right Search Method}{No method dominates everywhere}
  \small
  \renewcommand{\arraystretch}{1.35}
  \begin{center}
  \begin{tabularx}{\textwidth}{@{}>{\raggedright\arraybackslash}p{4.3cm} >{\bfseries\raggedright\arraybackslash}p{3.2cm} >{\raggedright\arraybackslash}X@{}}
    \toprule
    \rowcolor{PrimaryDark}
    \textcolor{white}{Situation} & \textcolor{white}{Method} & \textcolor{white}{Reason} \\
    \midrule
    Contingent planning, uncertain outcomes & AND-OR / AO-search & Must cover multiple continuations \\
    \rowcolor{PrimaryLight!25}
    Two-player game, manageable depth & Minimax & Optimal under perfect play \\
    Same setting, deeper trees & Alpha--beta & Minimax result, less search \\
    \rowcolor{PrimaryLight!25}
    Huge branching / weak heuristic & MCTS & Simulation focuses computation \\
    \bottomrule
  \end{tabularx}
  \end{center}
  \begin{center}\scriptsize
    Beyond board games: \keyword{autonomous driving} \textbullet\ \keyword{cybersecurity}
    \textbullet\ \keyword{robotics} \textbullet\ \keyword{auctions}.
  \end{center}
\end{frame}

\begin{frame}{Chapter Summary}{Key takeaways}
  \begin{columns}[T]
    \begin{column}{0.5\textwidth}
      \begin{block}{Modeling}
        \small
        \begin{itemize}
          \item Good decisions account for others' responses
          \item \keyword{AND-OR} trees capture contingencies
          \item Evaluation functions make big trees tractable
        \end{itemize}
      \end{block}
    \end{column}
    \begin{column}{0.5\textwidth}
      \begin{exampleblock}{Algorithms}
        \small
        \begin{itemize}
          \item \keyword{Minimax}: optimal in zero-sum games
          \item \keyword{Alpha--beta}: same answer, fewer nodes
          \item \keyword{MCTS}: simulation for huge trees
        \end{itemize}
      \end{exampleblock}
      \begin{alertblock}{Looking ahead}
        \small Ch.\,4: Propositional Logic and knowledge representation.
      \end{alertblock}
    \end{column}
  \end{columns}
\end{frame}

{\setbeamertemplate{footline}{}
\begin{frame}[plain,noframenumbering]
  \begin{tikzpicture}[remember picture,overlay]
    \fill[PrimaryDark](current page.north west) rectangle (current page.south east);
    \fill[PrimaryMid,opacity=0.35]
      ([shift={(0cm,-1.5cm)}]current page.north) circle(6.5cm);
  \end{tikzpicture}
  \vfill
  \begin{center}
    {\LARGE\textcolor{white}{\textbf{Thank you!}}}\\[0.6em]
    {\small\textcolor{PrimaryLight}{Questions \& Discussion}}\\[1.2em]
    {\footnotesize\textcolor{AccentYellow}{Next: Chapter 4 — Propositional Logic}}
  \end{center}
  \vfill
\end{frame}}

\end{document}
